\documentclass[10pt,conference]{IEEEtran}
\usepackage{cite}
\usepackage{amsmath,amssymb}
\usepackage{booktabs}
\usepackage{array}
\usepackage{url}
\usepackage[T1]{fontenc}
\renewcommand{\arraystretch}{1.08}
\begin{document}

\title{Why Can Only Some Cells Become Other Cells?\\
A Design-Only Framework for Testing Cellular Fate Accessibility and Reprogramming Constraints}

\author{Survey--Idea--ExperimentDesign--Author Research Workflow}

\maketitle

\begin{abstract}
Cells differ in their normal capacity to self-renew and produce other identities, but cellular fate is neither a binary property nor an absolutely irreversible one. Pluripotent stem cells, tissue-resident stem cells, progenitors, and differentiated cells occupy distinct regulatory and microenvironmental states. Induced pluripotency and direct lineage reprogramming demonstrate that selected differentiated cells can cross otherwise stable fate boundaries under specified perturbations, while incomplete conversion, residual source identity, genomic alteration, off-target lineages, uncontrolled proliferation, and host interaction limit interpretation and translation. This design-only paper reframes the question as a causal accessibility problem: for a declared source cell, target identity, niche, and perturbation, which regulatory and environmental gates permit a lineage-traced cell to acquire a stable, functional, and safe target state? We develop Fate Accessibility and Transition Evaluation - Lineage, Omics, Control, and Kinetics (FATE-LOCK), an auditable protocol that integrates source-provenance controls, time-resolved single-cell transcriptomic and chromatin assays, lineage tracing, perturbation/ablation/rescue logic, target-reference atlases, function and cue-withdrawal tests, genomic and lineage-safety gates, and independent replication. The framework requires a marker, trajectory, functional, and safety claim to be reported separately. It proposes no cell culture, reprogramming experiment, sequencing result, conversion efficiency, engraftment result, therapeutic product, or clinical outcome. Its conclusion is conditional: cellular fate is stabilized by coupled intrinsic and niche constraints, yet some boundaries can be experimentally traversed; credible evidence requires proving the route, identity, stability, function, and safety of each claimed transition.
\end{abstract}

\begin{IEEEkeywords}
stem cells, cell fate, differentiation, induced pluripotent stem cells, direct reprogramming, gene regulatory networks, stem-cell niche, single-cell multiomics, lineage tracing.
\end{IEEEkeywords}

\section{Introduction}
\label{sec:introduction}

The claim that only stem cells can become other cells captures an important observation about normal tissue organization, but it is too absolute as a scientific statement. Stem cells are distinguished by self-renewal and by an ability to generate one or more differentiated lineages under appropriate conditions. Embryonic stem cells and induced pluripotent stem cells can be pluripotent, whereas many tissue-resident adult stem cells act within a narrower, tissue-associated repertoire~\cite{nihstem}. Most mature cells maintain a characteristic identity rather than spontaneously becoming unrelated cells. However, induced pluripotency and direct lineage reprogramming show that selected mature cells can be pushed across fate boundaries using defined factors and conditions~\cite{takahashi2006,yu2007,graf2009}.

The central problem is thus not whether a cell carries a permanent label called ``transformative power.'' It is why some cellular states are accessible in a particular developmental, physiological, or engineered context and why others are strongly constrained. A source cell's capacity to transition depends on its lineage history, gene-regulatory network, chromatin state, transcription-factor availability, DNA methylation, replication and cell-cycle status, metabolism, stress response, and interactions with neighboring cells and extracellular matrix. These variables are coupled. Altering a transcription factor may change chromatin accessibility; altered accessibility can expose new regulatory interactions; niche signaling can favor survival of one subpopulation over another; and a new marker can arise in a stressed or partially reprogrammed cell without producing stable target function.

This distinction is consequential for regenerative medicine. A target-marker-positive population may still contain source-like cells, mixed-lineage intermediates, residual pluripotent cells, or proliferative cells that are unsuitable for a therapeutic product. An in-vitro transition may not persist after instructive cues are removed. Functional behavior may depend on a host niche unavailable in culture. A result in a lineage-traced model may therefore support a mechanistic transition claim while remaining far from tissue repair, cell-product qualification, or clinical safety.

This paper develops \emph{Fate Accessibility and Transition Evaluation - Lineage, Omics, Control, and Kinetics} (FATE-LOCK). FATE-LOCK is a design-only framework for a bounded question: given a specified source, target, microenvironment, and perturbation, can lineage-traced source cells enter a preregistered target-like state, retain it after cue withdrawal, exhibit target-relevant function, and pass the stated safety gates? The framework places source provenance, multimodal identity, trajectory, function, and safety on separate gates. It reports no culture, perturbation, sequencing, conversion, transplantation, disease-model, or clinical result.

\section{Survey: What Makes Cellular Fate Stable or Accessible?}
\label{sec:survey}

\subsection{Potency is graded and context-specific}

The native capacity of a cell is not captured by a single hierarchy alone. Early developmental cells, pluripotent cells, multipotent tissue stem cells, lineage-restricted progenitors, and mature differentiated cells differ in their usual developmental options. Yet their operational potential depends on assay conditions, developmental timing, species, tissue, and the criterion used for ``become.'' A cell can express a marker without acquiring full target function. It can transiently resemble a target transcriptome without stable self-maintenance. It can contribute to a tissue in a permissive assay without becoming broadly equivalent to the endogenous target population.

The NIH description of stem cells emphasizes both self-renewal and the distinction between pluripotent cells and adult stem cells associated with particular tissues~\cite{nihstem}. This provides a useful starting point, but it should not be read as an assertion that non-stem cells never change identity. Reprogramming experiments demonstrate an engineered alternative. Takahashi and Yamanaka showed that defined factors could induce pluripotent-like cells from mouse fibroblasts~\cite{takahashi2006}; Yu \emph{et al.}\ reported induced pluripotent stem-cell lines from human somatic cells~\cite{yu2007}. These studies do not eliminate the importance of native restriction. They show that restriction can sometimes be overcome by a specified molecular program in a specified model.

Table~\ref{tab:potency-boundary} summarizes the evidence boundary. The table avoids using words such as ``plastic'' or ``stemness'' as if they were universal assay-independent measurements. A claim becomes stronger as it moves from a marker to a stable, functional, and safe target identity.

\begin{table*}[!t]
\caption{Cell-state classes and evidence boundary. Potential is an assay- and context-dependent property, not a permanent universal label.}
\label{tab:potency-boundary}
\centering
\footnotesize
\begin{tabular}{@{}p{0.17\textwidth}p{0.26\textwidth}p{0.27\textwidth}p{0.22\textwidth}@{}}
\toprule
State class & Native role and accessible states & Evidence needed for a transition claim & Claim that remains unsupported \\
\midrule
Pluripotent stem cell & Self-renewal and differentiation toward adult-body lineages under directed conditions. & Identity, pluripotency controls, directed differentiation, target function, and safety assessment. & Every derivative is equivalent, mature, stable, or clinically usable. \\
Tissue-resident stem cell & Local maintenance and repair with tissue-associated lineage outputs. & Source provenance, niche context, lineage readout, and functional tissue-relevant assay. & Broad unrestricted potential outside the declared tissue and condition. \\
Lineage-restricted progenitor & Limited expansion and production of a subset of downstream identities. & Clonal/lineage proof, target comparison, and exclusion of rare contaminating precursors. & Self-renewal or pluripotency from a marker alone. \\
Differentiated cell under native conditions & Stable specialized function with usually restricted spontaneous fate change. & Time-resolved, lineage-traced evidence that a perturbation yields a stable target state. & Absolute irreversibility of identity. \\
Engineered reprogramming intermediate & Perturbation-dependent movement through a partial, direct, or pluripotent route. & Route, stability, source-memory, target function, proliferation, and off-target-lineage tests. & Complete target equivalence, tissue repair, or clinical safety. \\
\bottomrule
\end{tabular}
\end{table*}

\subsection{Gene regulatory networks and epigenetic barriers}

Cell identity can be modeled as a self-reinforcing regulatory configuration. Transcription factors activate or repress targets, recruit chromatin modifiers, and form feedback loops. Enhancers and promoters may be accessible in one state and inaccessible in another. DNA methylation, histone modification, nucleosome positioning, three-dimensional genome organization, and replication history influence which regulatory interactions are readily available. Young reviewed the coordinated transcriptional circuitry that controls embryonic stem-cell state~\cite{young2011}; this type of network helps explain why an identity can be stable even when the genome is shared across cell types.

For cell $i$, let $x_i(t)$ be a vector of regulatory, chromatin, protein, metabolic, and cell-cycle features; let $u_i(t)$ represent an engineered perturbation; let $n_i(t)$ represent niche inputs; and let $c_i(t)$ denote contextual history. A general state model is
\begin{align}
\frac{d x_i(t)}{dt}={}&f\!\left(x_i(t),u_i(t),n_i(t),c_i(t)\right) \nonumber\\
&+\eta_i(t),
\label{eq:state-dynamics}
\end{align}
where $f$ represents coupled biological dynamics and $\eta_i(t)$ represents unmodeled and stochastic variation. Equation~\eqref{eq:state-dynamics} is a conceptual model, not an assertion that all fate changes are deterministic or that each coordinate is measurable. It states why a transcription factor cannot be interpreted independently of cell state and environment.

One useful way to describe accessibility is as a target-specific probability rather than a permanent scalar. For target identities $j=1,\ldots,J$, define
\begin{equation}
\pi_{ij}(t)=\frac{\exp\{-\Phi_j(x_i(t),n_i(t),c_i(t))\}}{\sum_{\ell=1}^{J}\exp\{-\Phi_\ell(x_i(t),n_i(t),c_i(t))\}},
\label{eq:target-accessibility}
\end{equation}
where $\Phi_j$ is a preregistered model-dependent barrier or incompatibility score. Equation~\eqref{eq:target-accessibility} does not measure a literal physical energy barrier. It formalizes the idea that access to one target depends on both intrinsic and extrinsic state, and that changing $u_i(t)$ can alter the transition landscape without making every target equally reachable.

Reprogramming can be incomplete because the perturbation changes only part of this coupled system. A cell may lose selected source markers yet retain source chromatin or DNA-methylation memory. It may express a target transcription factor but fail to establish the enhancer network required for target function. It may survive only while exogenous cues are present. FATE-LOCK therefore requires source anti-markers, target anti-markers, multimodal comparisons, and stability after cue withdrawal rather than a single endpoint signature.

\subsection{The niche and tissue architecture}

The stem-cell niche is not merely a growth-factor container. It includes neighboring cells, extracellular matrix, physical confinement, oxygen and nutrient conditions, immune signals, neural or vascular input where relevant, and temporal cues. Scadden describes the niche as a regulatory neighborhood that affects stem-cell behavior~\cite{scadden2014}. The intestinal crypt illustrates how spatial organization and signaling support continuous tissue renewal in a defined compartment~\cite{clevers2013}. These examples caution against assuming that an in-vitro transcriptional resemblance is sufficient to recreate a tissue-resident fate.

Niche effects can act through several mechanisms. They may change transcription-factor activity, alter cell adhesion and mechanics, create survival selection, control quiescence or cell-cycle entry, and determine whether a newly induced state is maintained. They may also confound an experiment: a perturbation can enrich a rare target-like subpopulation instead of converting the declared source cell. FATE-LOCK therefore holds the niche explicit in the system boundary and uses lineage tracing and source-purity controls to distinguish selection from transition.

\subsection{Reprogramming, direct conversion, and the safety boundary}

Direct lineage reprogramming and induced pluripotency are powerful counterexamples to an absolute fate lock. Graf and Enver reviewed efforts to force lineage changes using regulatory factors~\cite{graf2009}. The mechanism can involve a pluripotent intermediate, a direct route, or partial conversion. These routes are not interchangeable. A direct route may avoid one intermediate but still have incomplete identity or off-target states; a pluripotent route may broaden differentiation options but introduces its own quality and residual-pluripotency concerns.

The transition from a research model to a clinical product adds independent requirements. Mandai \emph{et al.}\ described an autologous induced stem-cell-derived retinal-cell application in macular degeneration~\cite{mandai2017}. This work illustrates the importance of a controlled product and clinical protocol; it does not demonstrate that any reprogrammed cell product is safe or effective. Product identity, residual undifferentiated cells, genomic integrity, tumor-related risk, off-target lineages, manufacturing consistency, delivery, biodistribution, host interactions, and long-term follow-up cannot be inferred from a culture-level marker panel.

\section{Idea: FATE-LOCK as a Causal Accessibility Framework}
\label{sec:idea}

FATE-LOCK asks a bounded question: for a declared source population $s$, target identity $\tau$, culture or tissue niche $n$, and perturbation set $u$, which causal gates permit a lineage-traced source cell to reach a stable and functional target-like state without violating the stated safety boundary? The framework does not assume that a perfect target state is described by one transcriptomic centroid. It requires a reference definition that includes positive markers, anti-markers, chromatin features, morphology, function, stability, and the limitations of the model.

The central null hypothesis, $H_0$, is that an apparent conversion arises from source contamination, a rare pre-existing precursor, barcode collision, doublets, batch effects, selective survival, transient stress, or post hoc marker choice. Under $H_1$, a locked perturbation changes regulatory and niche gates so that lineage-traced source cells follow a reproducible trajectory and meet the predeclared identity, stability, and function criteria. Under $H_2$, selected markers arise without full function or stable target identity; the correct conclusion is partial or ambiguous reprogramming. Under $H_3$, access depends on source lineage, donor, niche, cell-cycle, genetic background, or perturbation timing, so a universal plasticity claim fails.

For a reference target profile $r_{jk}$ across feature domains $k=1,\ldots,K$, a transparent target similarity summary may be defined as
\begin{equation}
F_{ij}(t)=\sum_{k=1}^{K}w_k\,q_{ijk}(t),\qquad \sum_{k=1}^{K}w_k=1,
\label{eq:multimodal-score}
\end{equation}
where $q_{ijk}(t)$ is the calibrated, direction-consistent agreement of cell $i$ with target $j$ in domain $k$ and $w_k$ is locked before outcome review. Equation~\eqref{eq:multimodal-score} is a reporting summary, not a replacement for the individual domains. RNA similarity cannot compensate for absent target function, unrecognized source identity, or breached safety.

The framework has five principles. First, source identity and lineage provenance are measured before any target claim. Second, a target state is multimodal and includes negative evidence against source and off-target identities. Third, a transition is time-resolved, so an endpoint resemblance is not mistaken for a route. Fourth, function and cue-withdrawal persistence are tested separately from molecular resemblance. Fifth, safety is noncompensatory: a favorable target score does not offset genomic instability, uncontrolled proliferation, residual pluripotency where relevant, or inappropriate clinical extrapolation.

\begin{table*}[!t]
\caption{FATE-LOCK endpoint hierarchy. Each tier has a distinct failure mode and cannot be replaced by a favorable result from another tier.}
\label{tab:endpoint-hierarchy}
\centering
\footnotesize
\begin{tabular}{@{}p{0.17\textwidth}p{0.27\textwidth}p{0.27\textwidth}p{0.21\textwidth}@{}}
\toprule
Tier & Examples of locked measurements & Main inferential role & Prohibited shortcut \\
\midrule
Provenance & Source-cell audit, clone/barcode identity, sample-origin controls, doublet detection, and source spike-ins. & Excludes contamination and selection as explanations. & Calling a target-like contaminant a converted source cell. \\
State identity & RNA, chromatin, protein, morphology, positive markers, anti-markers, and reference-target comparison. & Tests whether the state resembles the declared target across modalities. & Treating one marker or one latent embedding as identity. \\
Trajectory & Multiple time points, lineage branches, perturbation timing, and route-specific intermediate definitions. & Tests whether a plausible source-to-target path is observed. & Inferring route from two endpoint snapshots. \\
Function and stability & Target-relevant assay, cue-withdrawal persistence, response behavior, and source-function loss where relevant. & Tests whether target-like state is durable and consequential in the model. & Equating transcriptional similarity with function. \\
Safety and translation boundary & Proliferation, residual pluripotency, genome integrity, off-target lineages, product purity, and host-interaction plan. & Limits the safety language that may be used. & Calling an in-vitro state transition a clinical cell therapy. \\
\bottomrule
\end{tabular}
\end{table*}

\section{ExperimentDesign: Lineage, Multiomics, Control, and Kinetics}
\label{sec:design}

\subsection{Scenario ladder and experimental boundary}

The protocol proceeds through S0 source audit, S1 baseline state atlas, S2 perturbation core, S3 stability/function, and S4 independent replication. Each scenario closes a different inferential dependency. S0 validates source purity and lineage provenance. S1 maps baseline heterogeneity without asserting causality. S2 estimates a perturbation contrast against matched controls. S3 tests whether induced features persist after cues are removed and have target-relevant consequences. S4 repeats a frozen protocol in an independent batch, site, or donor set. A later stage adds evidence but cannot erase a failed provenance or safety gate.

Table~\ref{tab:scenario-ladder} defines these scenarios. The boundary contains the source-cell class, donor or line provenance, target identity, medium and extracellular matrix, neighboring-cell or conditioned-medium inputs, oxygen/nutrient conditions, cell-cycle policy, perturbation vector, delivery method, time points, reference atlas, sequence and assay batches, identity/function/safety endpoints, and exclusion rules. All are frozen before the primary analysis. This is necessary because a source-to-target claim is only as interpretable as the cell population and environment that produced it.

\begin{table*}[!t]
\caption{FATE-LOCK scenario ladder. A completed lower rung never licenses a clinical or universal plasticity claim.}
\label{tab:scenario-ladder}
\centering
\footnotesize
\begin{tabular}{@{}p{0.15\textwidth}p{0.29\textwidth}p{0.28\textwidth}p{0.20\textwidth}@{}}
\toprule
Scenario & Required evidence and controls & Failure disqualifier & Permissible conclusion \\
\midrule
S0: source audit & Source identity markers and anti-markers, lineage or clonal provenance, contamination controls, doublet audit, and reference samples. & Source heterogeneity, rare precursor, or contamination cannot be excluded. & The stated source population is defined for the assay. \\
S1: state atlas & Time-resolved RNA/chromatin/protein baseline measurements, technical and biological replicates, batch and site controls. & State variation is confounded by untracked batch, stress, or cell-cycle distribution. & Baseline state variation is mapped under the stated conditions. \\
S2: perturbation core & Frozen factor/signaling/niche perturbation, matched controls, ablation, rescue, timing controls, and blinded assay labels. & Endpoint switching, source selection, or untracked delivery/culture variation. & A bounded perturbation-to-state-access contrast is estimated. \\
S3: stability and function & Target-relevant functional assay, cue withdrawal, persistence window, source-function assessment, and safety screening. & Target features are transient, nonfunctional, mixed-lineage, or safety criteria fail. & Target-like behavior is supported in the stated model. \\
S4: external replication & Locked protocol and analysis, independent operator/batch/site or donor set, null-result reporting. & Analysis is retuned after seeing replication data or effect reverses. & Reproducibility support is obtained for the replicated boundary. \\
\bottomrule
\end{tabular}
\end{table*}

The design does not prescribe a clinical-grade vector, a transplantation route, a therapeutic dose, or an intervention in people. Such choices require separate biosafety, manufacturing, toxicology, ethics, and regulatory development. The present protocol is intended to determine whether a mechanistic source-to-target claim is interpretable in its stated research model.

\subsection{Source provenance and perturbation logic}

At entry, each source cell or clone is assigned a versioned provenance record. This record includes collection or derivation history, passage or division history where relevant, genetic identity controls, source-marker and anti-marker panel, barcode or physical lineage-trace design, culture history, and batch. The reference target population is likewise versioned with its collection conditions, target markers, anti-markers, functional specification, and known heterogeneity. A target atlas does not become a circular definition of success: it is a comparator with stated limitations.

The perturbation core tests a preregistered factor, signaling, extracellular-matrix, neighboring-cell, or niche change against a matched control. It includes a perturbation-minus-key-factor or ablation arm and, where justified, a rescue arm. These controls separate necessity, sufficiency, timing, and survival effects. For example, if a perturbation increases target-marker frequency but an ablation leaves the effect unchanged, the initially proposed regulator is not required by the tested evidence. If only a pre-existing subpopulation survives, the result is selection rather than conversion.

For a binary target-occupancy state $T_i(t)$, the primary comparative estimate can be written as
\begin{equation}
\Delta_T(t)=\Pr[T_i(t)=1\mid A_i=1]-\Pr[T_i(t)=1\mid A_i=0],
\label{eq:occupancy-contrast}
\end{equation}
where $A_i$ indicates the frozen perturbation assignment. Equation~\eqref{eq:occupancy-contrast} must be stratified or modeled according to donor, clone, batch, source state, and other declared structure. It is not an estimate of conversion unless $T_i(t)=1$ includes the provenance, identity, trajectory, function, and safety criteria relevant to the claim.

Lineage evidence can be incorporated through a corrected target occupancy,
\begin{equation}
\widetilde{T}(t)=\frac{\sum_i L_i\,T_i(t)\,Q_i(t)}{\sum_i L_i\,Q_i(t)},
\label{eq:lineage-corrected-occupancy}
\end{equation}
where $L_i$ indicates verified source lineage and $Q_i(t)$ indicates that the cell passed the preregistered quality filter at time $t$. Equation~\eqref{eq:lineage-corrected-occupancy} makes clear that target occupancy is conditional on lineage and quality definitions. It does not correct for an invalid barcode, unmeasured contamination, or a post hoc quality filter; those risks require experimental and audit controls.

\subsection{Time-resolved multimodal measurement}

FATE-LOCK requires multiple time points because fate is a trajectory, not an endpoint image. RNA assays measure transcriptional state; chromatin assays estimate accessibility; protein assays measure a different layer; morphology and function provide orthogonal evidence. Cell-cycle and stress-state scores are collected so that a proliferating, dying, or acutely stressed cell is not misclassified as a developmental intermediate merely because it is distant from the source reference.

The observed feature vector may be decomposed as
\begin{equation}
Y_{ikt}=\theta_{ik}(t)+b_{k,\mathrm{batch}}+b_{k,\mathrm{site}}+\varepsilon_{ikt},
\label{eq:assay-decomposition}
\end{equation}
where $\theta_{ik}(t)$ is the underlying feature state, $b_{k,\mathrm{batch}}$ and $b_{k,\mathrm{site}}$ are calibrated batch and site effects, and $\varepsilon_{ikt}$ is residual variation. Equation~\eqref{eq:assay-decomposition} explains why batch, sample identity, and processing controls are part of causal interpretation rather than merely technical housekeeping.

Single-cell ordering can help formulate a putative route, as demonstrated by pseudotemporal methods~\cite{trapnell2014}. Such ordering does not itself prove ancestry or direction. Lineage tracing links state to fate more directly and can test whether cells sharing a transcriptional state have different later outputs~\cite{weinreb2020}. FATE-LOCK combines both: lineage tracing identifies the source relationship; time-resolved multiomics describes the state path; perturbation/ablation logic tests causal influence. No one layer replaces the others.

\subsection{Function, stability, and safety}

Target function must be specified before data collection. Its exact form depends on the target cell: an electrophysiological, secretory, contractile, barrier, metabolic, structural, or tissue-organization assay may be relevant. The protocol also tests stability after the instructive cue is withdrawn. A cell that satisfies a target marker only during forced expression or a specialized medium has a different status from a cell that autonomously maintains the state within the declared window. Both results can be scientifically informative, but their language must differ.

Safety is assessed relative to the research model and does not create a clinical safety certification. The protocol tracks uncontrolled proliferation, residual pluripotent-state markers when a pluripotent route is used, off-target lineages, genome-integrity signals appropriate to the method, and any model-specific adverse phenotype. A clinical product would require additional manufacturing consistency, product purity, delivery, biodistribution, tumorigenicity, immune interaction, dose, host-tissue integration, and long-term follow-up evidence. These requirements cannot be replaced by a high target-similarity score.

The noncompensatory gate is
\begin{equation}
G=\mathbb{I}_{\mathrm{provenance}}\mathbb{I}_{\mathrm{identity}}\mathbb{I}_{\mathrm{trajectory}}\mathbb{I}_{\mathrm{function}}\mathbb{I}_{\mathrm{stability}}\mathbb{I}_{\mathrm{safety}},
\label{eq:transition-gate}
\end{equation}
where every indicator must satisfy its preregistered criterion. Equation~\eqref{eq:transition-gate} ensures that $G=0$ if a target-like molecular pattern is found without lineage proof, function, stability, or safety. It prevents a numerical average from masking a failed gate.

\section{Systematics, Calibration, and Decision Rules}
\label{sec:systematics}

Before interpreting a conversion experiment, FATE-LOCK applies stress tests to the inference pipeline. Known source and target mixtures test whether the classifier detects contamination. Synthetic doublets test whether paired profiles appear as a false intermediate. Barcode collisions and sample swaps test provenance recovery. A known batch shift tests whether the normalization and reference alignment pipeline mistakes a technical transition for a biological one. Transient stress and cell-cycle shifts test whether trajectory tools falsely infer a fate branch. Null perturbation labels test whether the full analysis produces an effect when no causal assignment exists.

These tests do not simulate a positive biological result. They test whether the system is capable of rejecting incorrect explanations. A pipeline that calls conversion in a null-label set, cannot detect a source spike-in, or loses source lineage after integration cannot support a causal fate claim. Failed checks are reported, not silently removed. The required calibration material, reference atlases, quality controls, code version, and preprocessing parameters are fixed before evaluation.

For $m$ confirmatory hypotheses with ordered $p$-values $p_{(1)}\leq\cdots\leq p_{(m)}$, a preregistered family-wise step-down comparison can use
\begin{equation}
p_{(j)}\leq\frac{\alpha}{m-j+1}
\label{eq:holm-rule}
\end{equation}
at step $j$. Equation~\eqref{eq:holm-rule} makes the multiplicity family visible, but it cannot compensate for flexible source definitions, post hoc target-marker selection, or biased cell filtering. The confirmatory/exploratory distinction, target definitions, and feature-selection procedure must be locked in advance.

Table~\ref{tab:decision-gates} gives the progression logic. A passed gate supports a carefully limited statement. A failed later gate does not erase a valid earlier measurement finding, but it blocks stronger language. This makes it possible to learn from partial reprogramming, failed stability, or safety concerns without representing them as therapeutic success.

\begin{table*}[!t]
\caption{FATE-LOCK decision gates for a source-to-target transition claim. Each gate closes a distinct alternative explanation.}
\label{tab:decision-gates}
\centering
\footnotesize
\begin{tabular}{@{}p{0.15\textwidth}p{0.27\textwidth}p{0.28\textwidth}p{0.22\textwidth}@{}}
\toprule
Gate & Required evidence & Failure disqualifier & Permissible conclusion \\
\midrule
D0: source and target boundary & Locked source, target, niche, perturbation, time horizon, anti-markers, function, and safety definitions. & Identity words or endpoint definitions change after outcome inspection. & A testable cell-fate protocol exists. \\
D1: provenance & Source purity, barcode/clonal/physical lineage record, contamination and doublet controls, and sample identity audit. & Rare precursor, contaminant, or doublet explanation remains plausible. & The stated source is credible for the analysis. \\
D2: multimodal state and route & Predeclared RNA/chromatin/protein/morphology comparison, time points, batch controls, and ablation/rescue logic. & A single marker, embedding, or endpoint similarity substitutes for a route. & A bounded perturbation-associated target-like trajectory is supported. \\
D3: function and stability & Target-relevant function, persistence after cue withdrawal, source-function change, and complete null controls. & State is transient, nonfunctional, or mixed-lineage under the declared assay. & Stable target-like behavior is supported in the stated model. \\
D4: safety and replication & Genome/lineage/proliferation safety screen, independent replication, full failure and null reporting. & Safety gate fails, replication reverses, or analysis is retuned after replication data. & A reproducible research-model transition claim is supported. \\
\bottomrule
\end{tabular}
\end{table*}

\section{Mechanism Discrimination and Value of Information}
\label{sec:decision-ladder}

The purpose of a cell-fate experiment is not merely to maximize a target-marker fraction. It is to distinguish competing explanations for a restricted transition. A low target occupancy could arise because the source lacks a required regulator, because chromatin at target loci remains inaccessible, because the niche rejects an induced intermediate, because the perturbation is delivered at the wrong cell-cycle or developmental window, or because a target-like state is reached but cannot persist. These possibilities call for different experiments and lead to different conclusions. Treating all of them as a generic ``low efficiency'' obscures the mechanism that must be changed.

FATE-LOCK therefore treats every study as a sequence of decision-changing tests. If lineage tracing reveals that apparent target cells pre-existed in the input, then new transcriptional and chromatin assays do not answer the conversion question until source purity is repaired. If source cells reach a multimodal target-like state but lose it after cue withdrawal, then the open dependency is stabilization, not marker induction. If state and function are present in culture but disappear after exposure to a different matrix or neighboring-cell context, then the niche is part of the mechanism rather than an implementation detail. These conditional statements are more informative than a single score because they identify the next falsifiable hypothesis.

Table~\ref{tab:information-priorities} maps high-value uncertainties to discriminating measurements. The table is a planning aid, not a claim that the listed assays have been performed or that a particular result will occur. Its role is to prevent a resource-intensive sequencing campaign from being used when a less elaborate provenance or stability test would decide whether the biological question is even interpretable.

\begin{table*}[!t]
\caption{Mechanism-discrimination priorities. Each measurement is valuable only if it can change an experimental decision or close a stated alternative explanation.}
\label{tab:information-priorities}
\centering
\footnotesize
\begin{tabular}{@{}p{0.20\textwidth}p{0.26\textwidth}p{0.25\textwidth}p{0.21\textwidth}@{}}
\toprule
Uncertainty & Discriminating evidence & Decision changed by the evidence & What the evidence still does not prove \\
\midrule
Apparent target cells may be contaminants or rare precursors. & Lineage barcodes or clonal records, source spike-ins, target anti-markers, doublet controls, and independent sample identity checks. & Whether to test conversion mechanisms or first redefine/purify the source. & Complete target identity, function, or safety. \\
Target loci may be inaccessible despite factor expression. & Time-resolved chromatin and transcriptional measurements with factor ablation/rescue and matched delivery controls. & Whether a proposed regulator is required, sufficient in the model, or merely correlated. & A durable tissue contribution or clinical mechanism. \\
Niche conditions may select rather than convert cells. & Matrix, neighboring-cell, nutrient, oxygen, and cue-withdrawal comparisons with cell-count and death/proliferation tracking. & Whether to alter niche design, perturbation timing, or source selection logic. & That a culture phenotype will persist in a host tissue. \\
Target-like state may be transient or incomplete. & Functional assay, persistence after cue withdrawal, source-memory anti-markers, and longer declared observation window. & Whether the result remains a partial state, can proceed to stability work, or fails the target definition. & Cell-product qualification or human therapeutic safety. \\
Safety-relevant states may co-occur with target features. & Residual pluripotency, proliferation, off-target-lineage, and genome-integrity screens appropriate to the model. & Whether translational language is blocked and which safety mechanism requires study. & Absence of all long-term or in-vivo risk. \\
\bottomrule
\end{tabular}
\end{table*}

The value-of-information ordering also disciplines claims about reprogramming routes. A direct lineage transition and a pluripotent intermediate are not rhetorical alternatives; they lead to different quality questions. A route that appears direct still needs evidence against a transient intermediate, off-target branches, and residual source memory. A route that passes through a pluripotent state needs explicit residual-pluripotency, differentiation, proliferation, and lineage-purity checks. The appropriate experimental design follows the declared route and target, rather than treating the word ``reprogramming'' as a homogeneous mechanism.

Finally, FATE-LOCK makes the safety gate noncompensatory for an epistemic reason as well as a biomedical one. A high-dimensional state model may report a strong aggregate target score while a small subpopulation has an undesirable proliferative or off-target program. Averaging can hide the very cells that determine risk. The protocol therefore reports distributions, lineage branches, and failed safety criteria separately from average target similarity. It does not allow an attractive mean state to erase a discrete safety concern.

\section{Reproducibility, Translation, and Interpretation Safety}
\label{sec:reproducibility}

Every FATE-LOCK execution releases a versioned protocol, source and target definitions, lineage-trace schema, perturbation description, culture/niche parameters, sampling schedule, assay versions, batch identifiers, source and target reference atlases, quality filters, code, environment specifications, calibration controls, feature-selection rules, excluded-cell counts and reasons, null tests, synthetic-artifact tests, failed replicates, and data-access terms. Raw and processed data are shared subject to consent, governance, and biosafety restrictions. A run identifier changes whenever the source boundary, perturbation, target definition, reference atlas, quality filtering, or analysis plan changes.

Conclusions are labeled by evidence layer. An \emph{observation} is a measured marker, profile, barcode, or function readout. A \emph{state inference} is a comparison to a reference target under a measurement model. A \emph{causal transition inference} adds the perturbation and control assumptions. A \emph{translational claim} adds product, delivery, host, safety, and clinical assumptions. The framework forbids skipping directly from a cell-state observation to a patient-outcome statement.

The distinction is ethically important. People with severe disease may reasonably be interested in regenerative medicine, while experimental state conversion carries risks that cannot be inferred away. A research protocol must not promise repair, should disclose the difference between experimental cells and standard care, and should not use mechanistic excitement to reduce safety oversight. FATE-LOCK is intended to improve the quality of preclinical inference; it is not an authorization to use experimental reprogramming interventions in people.

\section{Threats to Validity and Research Priorities}
\label{sec:validity}

The dominant validity threat is source ambiguity. A rare precursor or contaminating target-like cell can make a culture appear to convert. Provenance controls, source spike-ins, anti-markers, and time-resolved lineage tracing address this threat, but their own error rates must be measured. A second threat is reference overfitting: a target atlas built from the same batch or selected after results can make any nearby cell look successful. Independent reference material and a locked feature set are necessary.

A third threat is incomplete reprogramming. Intermediate or stress states can share target markers, while residual source chromatin and function persist. The remedy is not simply more markers; it is an orthogonal combination of RNA, chromatin, protein, morphology, function, cue-withdrawal persistence, and negative evidence. A fourth threat is survival selection. If an intervention kills most source cells but spares a rare target-like population, increased target occupancy can be misread as conversion. Cell counts, lineage composition, death/proliferation measures, and source-linked trajectories must accompany occupancy estimates.

Time-scale mismatch is another concern. An early molecular transition may precede stable function by many divisions, or it may never mature. Conversely, a functional change may require a tissue niche not represented in culture. The protocol reports the observed horizon and model rather than extrapolating to lifelong tissue maintenance. The same logic applies to safety: an absence of an in-vitro signal does not demonstrate absence of a long-term or in-vivo risk.

Research priority should follow value of information. Measurements that change a decision gate deserve priority: lineage provenance; target anti-markers; independent reference controls; cue-withdrawal persistence; function; residual pluripotent and proliferative states; genomic integrity; and replication across source lineages, donors, and niches. A modest improvement in one target-marker frequency has limited value if it does not remove a contamination, stability, function, or safety uncertainty.

\section{Discussion}
\label{sec:discussion}

Why can only some cells become other cells? In normal biology, cells usually remain within stable identity programs because regulatory networks, chromatin configurations, developmental history, and niche signals reinforce their current state. Stem cells are specialized for self-renewal and developmental or tissue-maintenance outputs, so their native state makes certain transitions more accessible. But the premise must be qualified: engineered reprogramming and direct conversion show that selected differentiated cells can cross some boundaries. The meaningful scientific question is which boundaries, routes, and conditions are accessible, and whether the resulting state is complete, stable, functional, and safe.

FATE-LOCK converts this question into an auditable sequence. A researcher begins with a source and target definition, verifies lineage, maps baseline heterogeneity, tests a locked perturbation against controls, measures a time-resolved multimodal route, tests function and cue-withdrawal stability, screens safety, and replicates the result. Each stage eliminates a different alternative explanation. A target-like marker alone may be interesting, but it cannot satisfy the final gate.

The framework also clarifies how to communicate reprogramming research. A successful S2 transition can support a molecular or mechanistic result in a defined model. A successful S3 result can support target-like behavior in that model. An S4 replication can improve confidence in the bounded research claim. None of these automatically establishes universal cellular plasticity, organ regeneration, or human therapy. This is not a limitation of ambition; it is the evidence structure needed to make genuine progress distinguishable from a persuasive but incomplete cell-state image.

\section{Conclusion}
\label{sec:conclusion}

This four-stage research workflow concludes that cell-fate restriction is real but conditional, not an absolute prohibition on identity change. Coupled gene-regulatory, epigenetic, lineage-history, and niche mechanisms stabilize cellular states, while selected perturbations can make some transitions accessible. A credible claim that a source cell became another cell requires proof of source provenance, multimodal target identity, time-resolved route, target-relevant function, stability after cue withdrawal, and a stated safety boundary. FATE-LOCK provides this design-only benchmark. It reports no performed reprogramming or therapeutic result, and it does not claim that an engineered state transition is clinically safe.

\appendices
\section{Minimal FATE-LOCK Run Manifest}
\label{app:manifest}

Each future run documents: source and target definitions; source history and provenance; donor/clone and genetic-identity controls; lineage-trace design; culture, matrix, niche, nutrient, oxygen, and timing conditions; perturbation and matched controls; reference atlas versions; assay and batch metadata; positive and anti-markers; quality filters; function and cue-withdrawal assays; proliferation, residual-pluripotency, off-target-lineage, and genome-integrity screens; barcode/doublet/sample-swap controls; all exclusions; all null and synthetic-artifact tests; replication plan; code and environment versions; and a statement that the model cannot certify a clinical cell product or human therapeutic outcome.

\section{Protocol Checklist}
\label{app:checklist}

Before an experiment, freeze source, target, niche, perturbation, route, identity, function, stability, and safety definitions. Before analysis, verify source provenance, sample identity, batch controls, duplicate measurements, reference atlas independence, and assay-quality thresholds. During interpretation, report target markers and anti-markers, every cell-state component, source/target cell counts, proliferation/death context, failed replicates, excluded cells, and null tests. Before communication, label the conclusion as a measurement observation, state inference, causal transition inference, or translational extrapolation; do not use clinical or universal-plasticity language for a finite research-model result.

\begin{thebibliography}{10}

\bibitem{nihstem}
National Institutes of Health, ``Stem Cell Basics,'' NIH Stem Cell Information, accessed Sep. 2, 2026. [Online]. Available: \url{https://stemcells.nih.gov/info/basics/stc-basics}

\bibitem{takahashi2006}
K. Takahashi and S. Yamanaka, ``Induction of pluripotent stem cells from mouse embryonic fibroblasts by defined factors,'' \emph{Cell}, vol. 126, no. 4, pp. 663--676, 2006, doi: 10.1016/j.cell.2006.07.024.

\bibitem{yu2007}
J. Yu \emph{et al.}, ``Induced pluripotent stem cell lines derived from human somatic cells,'' \emph{Science}, vol. 318, no. 5858, pp. 1917--1920, 2007, doi: 10.1126/science.1151526.

\bibitem{graf2009}
T. Graf and T. Enver, ``Forcing cells to change lineages,'' \emph{Nature}, vol. 462, pp. 587--594, 2009, doi: 10.1038/nature08533.

\bibitem{young2011}
R. A. Young, ``Control of the embryonic stem cell state,'' \emph{Cell}, vol. 144, no. 6, pp. 940--954, 2011, doi: 10.1016/j.cell.2011.01.032.

\bibitem{scadden2014}
D. T. Scadden, ``Nice neighborhood: Emerging concepts of the stem cell niche,'' \emph{Cell}, vol. 157, no. 1, pp. 41--50, 2014, doi: 10.1016/j.cell.2014.02.013.

\bibitem{clevers2013}
H. Clevers, ``The intestinal crypt, a prototype stem cell compartment,'' \emph{Cell}, vol. 154, no. 2, pp. 274--284, 2013, doi: 10.1016/j.cell.2013.07.004.

\bibitem{trapnell2014}
C. Trapnell \emph{et al.}, ``The dynamics and regulators of cell fate decisions are revealed by pseudotemporal ordering of single cells,'' \emph{Nature Biotechnology}, vol. 32, no. 4, pp. 381--386, 2014, doi: 10.1038/nbt.2859.

\bibitem{weinreb2020}
C. Weinreb \emph{et al.}, ``Lineage tracing on transcriptional landscapes links state to fate during differentiation,'' \emph{Science}, vol. 367, no. 6472, Art. eaaw3381, 2020, doi: 10.1126/science.aaw3381.

\bibitem{mandai2017}
M. Mandai \emph{et al.}, ``Autologous induced stem-cell-derived retinal cells for macular degeneration,'' \emph{New England Journal of Medicine}, vol. 376, no. 11, pp. 1038--1046, 2017, doi: 10.1056/NEJMoa1608368.

\end{thebibliography}

\end{document}
